% (c) Nikita Lisitsa, lisyarus@gmail.com, 2025

\documentclass[10pt]{beamer}

\usepackage[T2A]{fontenc}
\usepackage[russian]{babel}
\usepackage{minted}

\usepackage[most]{tcolorbox}
\tcbuselibrary{minted}

\usepackage{graphicx}
\graphicspath{ {./images/} }

\usepackage{adjustbox}

\usepackage{color}
\usepackage{soul}

\usepackage{hyperref}

\usetheme{metropolis}

\definecolor{red}{rgb}{1,0,0}
\definecolor{green}{rgb}{0,0.5,0}
\definecolor{codebg}{RGB}{29,35,49}
\setminted{fontsize=\footnotesize}

\definecolor{blue}{rgb}{0,0,1}
\definecolor{red}{rgb}{1,0,0}

\makeatletter
\newcommand{\slideimage}[1]{
  \begin{figure}
    \begin{adjustbox}{width=\textwidth, totalheight=\textheight-2\baselineskip-2\baselineskip,keepaspectratio}
      \includegraphics{#1}
    \end{adjustbox}
  \end{figure}
}
\makeatother

\usemintedstyle{tango}

\title{Компьютерная графика}
\subtitle{Лекция 2: Асинхронность, растеризация, графический конвейер, шейдеры, интерполяция, аффинные преобразования}
\date{2025}

\setbeamertemplate{footline}[frame number]

\begin{document}

\frame{\titlepage}

\begin{frame}
\frametitle{Асинхронность GPU}
\begin{itemize}
\item GPU -- отдельное периферийное устройство
\pause
\begin{itemize}
\item К нему нет прямого доступа
\pause
\item Одновременно используется несколькими процессами
\pause
\item Одновременно выполняет несколько задач
\end{itemize}
\pause
\item Все команды оно выполняет \textit{асинхронно}
\end{itemize}
\end{frame}

\begin{frame}
\frametitle{Асинхронность GPU и SIMD}
\pause
\begin{itemize}
\item SIMD -- Single Instruction, Multiple Data
\pause
\item GPU -- как SIMD-расширения процессоров (SSE, AVX), но на более крупном масштабе: сотни/тысячи одновременных вычислений
\pause
\item Для производительности вычислительные ядра организованы в группы (warps/wavefronts/SM) по 32/64 элемента
\pause
\item Все ядра одного wavefront'а выполняют одну и ту же операцию (но над разными данными)
\pause
\item Если мы рисуем только один треугольник, под него всё равно выделяется целый wavefront и большинство его ядер делают бесполезную работу
\pause
\item Максимизация загрузки wavefront'ов -- важная часть низкоуровневой оптимизации GPU
\end{itemize}
\end{frame}

\begin{frame}
\frametitle{Асинхронность GPU}
\begin{itemize}
\item Старые графические API (напр. OpenGL) прятали это в реализацию драйвера
\pause
\item Современные API (напр. WebGPU) явно отражают это в своём интерфейсе
\pause
\item {\color{red}\texttt{-}} Нет операции ``нарисовать объект''
\item {\color{green}\texttt{+}} Есть операция ``записать последовательность команд и отправить в очередь на исполнение''
\end{itemize}
\end{frame}

\begin{frame}
\frametitle{Queue, command encoder, command buffer}
\begin{itemize}
\item \mintinline{rust}|CommandEncoder| -- объект, позволяющий записать набор команд для исполнения на GPU
\pause
\item \mintinline{rust}|encoder.finish() -> CommandBuffer| -- записанный буфер с командами в понятном драйверу/GPU виде
\pause
\item \mintinline{rust}|queue.submit(buffers)| -- отправить команды на исполнение (принимает массив \mintinline{rust}|CommandBuffer|'ов)
\end{itemize}
\end{frame}

\begin{frame}
\frametitle{Растеризация}
\begin{itemize}
\item Растеризация -- превращение геометрического примитива (точки, линии, треугольника, прямоугольника, круга, и т.д.) в набор соответствующих ему пикселей на экране/изображении
\pause
\item Превращение векторных данных в растровые
\pause
\item За нас её делает GPU!
\pause
\item Некоторые современные графические движки (Unreal 5 Nanite) делают растеризацию сами с помощью compute шейдеров
\end{itemize}
\end{frame}

\begin{frame}
\frametitle{Почему растеризация?}
\begin{itemize}
\item Исторически было очень много разных способов рисовать трёхмерные сцены; см. \href{https://dl.acm.org/doi/pdf/10.1145/356625.356626}{\texttt{Sutherland et al - A Characterization of Ten Hidden-Surface Algorithms}}
\pause
\item В современности есть два основных алгоритма:
\pause
\begin{itemize}
\item \textbf{Растеризация}: какие пиксели занимает конкретный объект
\pause
\item \textbf{Трассировка лучей}: какие объекты видит конкретный пиксель
\end{itemize}
\end{itemize}
\end{frame}

\begin{frame}
\frametitle{Почему растеризация?}
\begin{itemize}
\item Особенности растеризации:
\begin{itemize}
\item {\color{green}\texttt{+}} Быстрее (зависит от сцены)
\pause
\item {\color{green}\texttt{+}} Легче параллелится
\pause
\item {\color{green}\texttt{+}} Лучше ложится на железо
\pause
\item {\color{green}\texttt{+}} Лучше работает с памятью (data locality)
\pause
\item {\color{red}\texttt{-}} Сложнее добиться фотореалистичной графики
\end{itemize}
\pause
\item Особенности трассировки лучей:
\pause
\begin{itemize}
\item Всё ровно наоборот :)
\end{itemize}
\end{itemize}
\end{frame}

\begin{frame}[fragile]
\frametitle{Растеризация: точка}
\begin{itemize}
\item Как растеризовать точку \begin{math}(x, y)\end{math}?
\pause
\usemintedstyle{lightbulb}
\begin{tcblisting}{listing engine=minted, minted language=cpp, minted options={fontsize=\scriptsize}, listing only, colback=codebg, colframe=codebg, rounded corners}
set_pixel(floor(x), floor(y), color);
\end{tcblisting}
\pause
\usemintedstyle{tango}
\item В WebGPU: \mintinline{cpp}|point-list|
\end{itemize}
\end{frame}

\begin{frame}
\frametitle{Растеризация: точка}
\slideimage{raster-point.png}
\end{frame}

\begin{frame}[fragile]
\frametitle{Растеризация: линия}
\begin{itemize}
\item Как растеризовать линию \begin{math}(x_1, y_1) \dots (x_2, y_2)\end{math}?
\pause
\item Алгоритм Брезенхэма
\pause
\item Есть вариация алгоритма для рисования окружностей
\pause
\usemintedstyle{tango}
\item В WebGPU: \mintinline{cpp}|line-list|, \mintinline{cpp}|line-strip|
\end{itemize}
\end{frame}

\begin{frame}
\frametitle{Растеризация: линия}
\slideimage{raster-line.png}
\end{frame}

\begin{frame}
\frametitle{Растеризация: окружность}
\slideimage{raster-circle.png}
\end{frame}

\begin{frame}[fragile]
\frametitle{Растеризация: прямоугольник}
\begin{itemize}
\item Как растеризовать прямоугольник \begin{math}[x_1\dots x_2]\times[y_1\dots y_2]\end{math}?
\end{itemize}
\pause
\usemintedstyle{lightbulb}
\begin{tcblisting}{listing engine=minted, minted language=cpp, minted options={fontsize=\scriptsize}, listing only, colback=codebg, colframe=codebg, rounded corners}
for (int x = round(x_1); x < round(x_2); ++x) {
    for (int y = round(y_1); y < round(y_2); ++y) {
        set_pixel(x, y, color);
    }
}
\end{tcblisting}
\end{frame}

\begin{frame}
\frametitle{Растеризация: прямоугольник}
\slideimage{raster-rect.png}
\end{frame}

\begin{frame}[fragile]
\frametitle{Растеризация: треугольник}
\begin{itemize}
\item Как растеризовать треугольник с вершинами \begin{math}(x_1, y_1), (x_2, y_2), (x_3, y_3)\end{math}?
\pause
\item Растеризуем ограничивающий прямоугольник, проверяя пиксели на вхождение в треугольник
\pause
\usemintedstyle{lightbulb}
\begin{tcblisting}{listing engine=minted, minted language=cpp, minted options={fontsize=\scriptsize}, listing only, colback=codebg, colframe=codebg, rounded corners}
int xmin = min(floor(x_1), floor(x_2), floor(x_3));
int xmax = max( ceil(x_1),  ceil(x_2),  ceil(x_3));

int ymin = min(floor(y_1), floor(y_2), floor(y_3));
int ymax = max( ceil(y_1),  ceil(y_2),  ceil(y_3));

for (int x = xmin; x <= xmax; ++x) {
    for (int y = ymin; y <= ymax; ++y) {
        if (inside_triangle(x, y, ...)) {
            set_pixel(x, y, color);
        }
    }
}
\end{tcblisting}
\pause
\usemintedstyle{tango}
\item В WebGPU: \mintinline{cpp}|triangle-list|, \mintinline{cpp}|triangle-strip|
\end{itemize}
\end{frame}

\begin{frame}
\frametitle{Растеризация: треугольник}
\slideimage{raster-triangle.png}
\end{frame}

\begin{frame}[fragile]
\frametitle{Растеризация: круг}
\begin{itemize}
\item Как растеризовать круг с центром \begin{math}(x_0, y_0)\end{math} и радиусом \begin{math}R\end{math}?
\pause
\item Растеризуем ограничивающий прямоугольник, проверяя пиксели на вхождение в круг
\end{itemize}
\pause
\usemintedstyle{lightbulb}
\begin{tcblisting}{listing engine=minted, minted language=cpp, minted options={fontsize=\scriptsize}, listing only, colback=codebg, colframe=codebg, rounded corners}
int xmin = floor(x_0 - R);
int xmax =  ceil(x_0 + R);

int ymin = floor(y_0 - R);
int ymax =  ceil(y_0 + R);

for (int x = xmin; x <= xmax; ++x) {
    for (int y = ymin; y <= ymax; ++y) {
        if (sqr(x - x_0) + sqr(y - y_0) <= sqr(R))
            set_pixel(x, y, color);
    }
}
\end{tcblisting}
\end{frame}

\begin{frame}
\frametitle{Растеризация: круг}
\slideimage{raster-disk.png}
\end{frame}

\begin{frame}<1>[fragile,label=rasterization]
\frametitle{Правила растеризации}
\begin{itemize}
\item Пиксель растеризуется, если \textit{центр пикселя} содержится в треугольнике
\pause
\item Если у двух треугольников есть общее ребро (и они не пересекаются внутренностями), то
\begin{itemize}
\item Каждый пиксель будет принадлежать ровно одному треугольнику, т.е. не будет наложения
\item Ни один пиксель общего ребра не будет пропущен, т.е. не будет ``дырок''
\end{itemize}
\pause
\item Подробнее: \href{https://en.wikibooks.org/wiki/GLSL_Programming/Rasterization}{\nolinkurl{en.wikibooks.org/wiki/GLSL\_Programming/Rasterization}}
\end{itemize}
\end{frame}

\begin{frame}
\frametitle{Правила растеризации}
\slideimage{pixel-covered.png}
\end{frame}

\againframe<1-2>{rasterization}

\begin{frame}
\frametitle{Правила растеризации}
\:
\slideimage{triangle-rasterization.png}
\end{frame}

\begin{frame}
\frametitle{Правила растеризации}
Не будет наложения пикселей:
\slideimage{triangle-rasterization-overlap.png}
\end{frame}

\begin{frame}
\frametitle{Правила растеризации}
Не будет дырок:
\slideimage{triangle-rasterization-hole.png}
\end{frame}

\againframe<2-3>{rasterization}

\begin{frame}[fragile]
\frametitle{Правила растеризации}
\begin{itemize}
\item В современных графических API есть только три основных примитива для растеризации:
\pause
\begin{itemize}
\item \textbf{Точки}: \mintinline{cpp}|point-list|
\pause
\item \textbf{Линии}: \mintinline{cpp}|line-list|, \mintinline{cpp}|list-strip|
\pause
\item \textbf{Треугольники}: \mintinline{cpp}|triangle-list|, \mintinline{cpp}|triangle-strip|
\end{itemize}
\pause
\item Точки и линии обычно используются для дебажного рендеринга или вместе с геометрическими шейдерами, \textit{главный примитив -- \underline{треугольник}}
\end{itemize}
\end{frame}

\begin{frame}
\frametitle{О треугольниках}
\begin{itemize}
\item Почему основным примитивом рисования стал \textit{треугольник}?
\pause
\item Для более сложных геометрических фигур (круг, многоугольник, и т.д.):
\pause
\begin{itemize}
\item {\color{red}\texttt{-}} Сложнее алгоритм разтеризации (особенно с учётом перспективной проекции)
\pause
\item {\color{red}\texttt{-}} Сложнее задавать и интерполировать атрибуты (накладывать цвет и текстуру, вычислять освещение)
\end{itemize}
\end{itemize}
\end{frame}

\begin{frame}
\frametitle{О треугольниках}
Плюсы треугольника:
\pause
\begin{itemize}
\item {\color{green}\texttt{+}} Образ под действием перспективной проекции -- тоже треугольник
\pause
\item {\color{green}\texttt{+}} Есть единственный разумный способ интерполяции (линейная, с барицентрическими координатами)
\pause
\item {\color{green}\texttt{+}} Позволяет рисовать спрайты (прямоугольник -- два треугольника)
\pause
\item {\color{green}\texttt{+}} Позволяет рисовать многоугольники (посредством триангуляции)
\pause
\item {\color{green}\texttt{+}} Позволяет рисовать линии (превращая их в тонкие многоугольники)
\pause
\item {\color{green}\texttt{+}} Позволяет рисовать более сложные фигуры (аппроксимируя)
\end{itemize}
\end{frame}

\begin{frame}[fragile]
\frametitle{Графический конвейер (graphics pipeline)}
\begin{itemize}
\item Последовательность операций на GPU, превращающих входные данные в картинку на экране
\pause
\item Что происходит при выполнении \mintinline{cpp}|render_pass.draw()| от чтения данных до появления картинки на экране
\pause
\item Почти не отличается для разных графических API
\end{itemize}
\end{frame}

\begin{frame}[fragile]
\frametitle{Графический конвейер (graphics pipeline)}
Основные части конвейера:
\begin{itemize}
\item Входной поток вершин (vertex stream)
\pause
\begin{itemize}
\item Позже мы узнаем, как его задавать
\end{itemize}
\pause
\item Вершинный шейдер: обрабатывает вершины по одной
\pause
\item Сборка примитивов (primitive assembly) из отдельных вершин
\pause
\item Преобразование в оконную систему координат (viewport transform)
\pause
\item Растеризация примитивов
\pause
\item Пиксельный (фрагментный) шейдер: обрабатывает пиксели по одному
\end{itemize}
\end{frame}

\begin{frame}[fragile]
\frametitle{Графический конвейер (graphics pipeline)}
\begin{itemize}
\item Мы пропустили много важных частей конвейера
\item Будем их по чуть-чуть добавлять в течение курса
\item \href{https://www.khronos.org/opengl/wiki/Rendering_Pipeline_Overview}{\nolinkurl{khronos.org/opengl/wiki/Rendering_Pipeline_Overview}}
\end{itemize}
\end{frame}

\begin{frame}
\frametitle{Графический конвейер (graphics pipeline)}
\begin{itemize}
\item Конвейер -- удобная абстракция, внутри GPU всё устроено сложнее
\pause
\item Каждая часть конвейера выполняется сразу на большом объёме входных данных параллельно
\pause
\item Разные части конвейера могут выполняться одновременно (на разных данных)
\pause
\item Разные вызовы команд рисования могут обрабатываться параллельно
\end{itemize}
\end{frame}

\begin{frame}[fragile]
\frametitle{Вершинный (vertex) шейдер}
\begin{itemize}
\pause
\item Входные данные:
\pause
\begin{itemize}
\item Аттрибуты вершин (мы позже узнаем, как их задавать) -- свои значения для каждой вершины (позиция, цвет, нормаль)
\pause
\item Immediates/push constants -- маленький (128 байт) набор данных, вшитый в command buffer
\pause
\item Встроенные переменные, например \mintinline{wgsl}|@builtin(vertex_index)| -- полный список смотри в \href{https://gpuweb.github.io/gpuweb/wgsl/#builtin-inputs-outputs}{\nolinkurl{документации}}
\pause
\item Uniform'ы, текстуры, буферы -- об этом позже :)
\end{itemize}
\pause
\item Выходные данные:
\begin{itemize}
\item \mintinline{wgsl}|@builtin(position) position: vec4f|
\pause
\item Переменные, интерполированное значение которых попадёт во фрагментный (пиксельный) шейдер: \mintinline{wgsl}|@location(0) color: vec4f|
\end{itemize}
\end{itemize}
\end{frame}

\begin{frame}<1>[label=primitive_assembly]
\frametitle{Сборка примитивов (primitive assembly)}
\begin{itemize}
\item Превращает набор независимых вершин в \textit{притимивы} растеризации: точки/линии/треугольники
\pause
\item Правила сборки зависят от выбранного типа примитива: \mintinline{cpp}|point-list|, \mintinline{cpp}|line-list|, \mintinline{cpp}|line-strip|, \mintinline{cpp}|triangle-list|, \mintinline{cpp}|triangle-strip|
\end{itemize}
\end{frame}

\begin{frame}
\frametitle{Сборка примитивов (primitive assembly)}
\slideimage{primitives.png}
\end{frame}

\againframe<1->{primitive_assembly}

\begin{frame}[fragile]
\frametitle{Viewport transform}
\begin{itemize}
\item Преобразование в оконную (пиксельную) систему координат
\item \begin{math}X: [-1, 1] \rightarrow [X_{min}, X_{max}]\end{math}
\item \begin{math}Y: [-1, 1] \rightarrow [Y_{max}, Y_{min}]\end{math} (-1 внизу, 1 вверху)
\pause
\item Настроить: \mintinline{cpp}|render_pass.set_viewport()|
\pause
\item По умолчанию -- весь экран (точнее, весь color target)
\pause
\item Вне прямоугольника \begin{math}[X_{min}, X_{max}]\times[Y_{min}, Y_{max}]\end{math} \textit{ничего не нарисуется}: GPU обрежет все примитивы по этим границам
\end{itemize}
\end{frame}

\begin{frame}[fragile]
\frametitle{Фрагментный (пиксельный, fragment) шейдер}
\begin{itemize}
\pause
\item Входные данные:
\pause
\begin{itemize}
\item Проинтерполированные значения из вершинного шейдера: \mintinline{wgsl}|@location(0) color: vec4f|
\pause
\item Immediates/push constants
\pause
\item Встроенные переменные, например \mintinline{wgsl}|@builtin(position)|
\pause
\item Uniform'ы, текстуры, буферы
\end{itemize}
\pause
\item Выходные данные:
\begin{itemize}
\item \mintinline{wgsl}|@location(0) color: vec4f| -- выходной цвет в формате RGBA, компоненты в диапазоне \([0\dots 1]\)
\pause
\item Может быть несколько, об этом поговорим позже
\end{itemize}
\end{itemize}
\end{frame}

\begin{frame}[fragile]
\frametitle{Render pass encoder}
\begin{itemize}
\item При рисовании нам нужно сказать, \textbf{куда} рисовать, \textbf{что} рисовать, и \textbf{как} рисовать
\pause
\item Куда (в какую текстуру) рисовать настраивается при создании \mintinline{rust}|RenderPassEncoder|'а
\pause
\item Что рисовать настраивается через методы \mintinline{rust}|RenderPassEncoder|'а (\mintinline{rust}|set_vertex_buffer| и т.п.)
\pause
\item Как рисовать настраивается через \mintinline{rust}|render_pass_encoder.set_pipeline()|
\end{itemize}
\end{frame}

\begin{frame}[fragile]
\frametitle{Render pipeline}
\begin{itemize}
\item Часть настроек конвейера -- тяжёлые и сложные, с большим объёмом валидации: настройки атрибутов вершин, используемые шейдеры, тип примитива, и т.п.
\pause
\item Часть настроек -- лёгкие, часто описываются парой чисел: viewport, scissor, blend constant, stencil reference
\pause
\item Тяжёлые настройки хранит особый объект -- \mintinline{rust}|RenderPipeline|
\pause
\item Лёгкие настройки можно менять на лету через \mintinline{rust}|RenderPassEncoder|
\end{itemize}
\end{frame}

\begin{frame}[fragile]
\frametitle{Языки описания шейдеров}
\begin{itemize}
\item \textbf{OpenGL, Vulkan}: \textbf{GLSL} (GL Shading Language)
\pause
\item \textbf{DirectX}: \textbf{HLSL} (High-Level Shading Language)
\pause
\item \textbf{DirectX} (до 2012): \textbf{Cg} (C for Graphics), deprecated
\pause
\item \textbf{WebGPU}: \textbf{\alert{\underline{WGSL}}} (WebGpu Shading Language)
\pause
\item \textbf{Slang} -- язык, не привязанный к API
\pause
\item И другие: \href{https://en.wikipedia.org/wiki/Shading_language}{\nolinkurl{en.wikipedia.org/wiki/Shading\_language}}
\end{itemize}
\end{frame}

\begin{frame}[fragile]
\frametitle{Язык описания шейдеров WGSL}
\begin{itemize}
\item Похож на Rust
\pause
\item Типы данных:
\pause
\begin{itemize}
\item Скалярные: \mintinline{wgsl}|bool|, \mintinline{wgsl}|i32|, \mintinline{wgsl}|u32|, \mintinline{wgsl}|f32|
\pause
\item Векторные: \mintinline{wgsl}|vec2<T>|, \mintinline{wgsl}|vec3<T>|, \mintinline{wgsl}|vec4<T>|
\pause
\item Векторные алиасы: \mintinline{wgsl}|vec2i|, \mintinline{wgsl}|vec3f|, ...
\pause
\item Матричные: \mintinline{wgsl}|mat2x4<T>|, \mintinline{wgsl}|mat4x4f|, ...
\pause
\item \textbf{N.B.:} \mintinline{wgsl}|mat2x4| это матрица \(4\times 2\), т.е. с 4 строками и 2 столбцами!
\end{itemize}
\end{itemize}
\end{frame}

\begin{frame}[fragile]
\frametitle{Язык описания шейдеров WGSL}
\begin{itemize}
\item Конструкторы векторов:
\pause
\begin{itemize}
\item \mintinline{wgsl}|vec3f(x, y, z)|
\pause
\item \mintinline{wgsl}|vec3f(x) == vec3f(x, x, x)|
\pause
\item \mintinline{wgsl}|vec3f(vec2f(x, y), z) == vec3f(x, vec2f(y, z))|
\end{itemize}
\pause
\item Подробнее: \href{hhttps://google.github.io/tour-of-wgsl/types}{\nolinkurl{google.github.io/tour-of-wgsl/types}}
\end{itemize}
\end{frame}

\begin{frame}[fragile]
\frametitle{Язык описания шейдеров WGSL}
\begin{itemize}
\item Есть стандартные операции: \mintinline{wgsl}|+|, \mintinline{wgsl}|-|, \mintinline{wgsl}|*|, \mintinline{wgsl}|/|, \mintinline{wgsl}|%|, \mintinline{wgsl}|<|, \mintinline{wgsl}|==|, ...
\pause
\begin{itemize}
\item Для векторов и матриц применяются покомпонентно
\pause
\item Исключение: умножение двух матриц это умножение в смысле алгебры
\pause
\item Исключение: умножение матрицы на вектор слева это умножение в смысле алгебры
\pause
\item Исключение: умножение матрицы на вектор справа это умножение в смысле алгебры на транспонированную матрицу
\end{itemize}
\pause
\item Доступ к координатам векторов: \mintinline{wgsl}|a.x = b.y|
\pause
\item Swizzling: \mintinline{wgsl}|a.xxzy == vec4f(a.x, a.x, a.z, a.y)|
\pause
\item Доступ к элементам матриц: \mintinline{wgsl}|m[column][row]|
\pause
\item Есть полезные математические функции: \mintinline{wgsl}|pow|, \mintinline{wgsl}|sin|, \mintinline{wgsl}|cos|, ...
\pause
\begin{itemize}
\item Для векторов и матриц применяются покомпонентно
\end{itemize}
\pause
\item Есть математические функции векторной алгебры: \mintinline{wgsl}|dot|, \mintinline{wgsl}|cross|, \mintinline{wgsl}|length|, \mintinline{wgsl}|normalize|, ...
\pause
\item Линейная и нелинейная интерполяция: \mintinline{wgsl}|mix|, \mintinline{wgsl}|smoothstep|
\end{itemize}
\end{frame}

\begin{frame}[fragile]
\frametitle{Язык описания шейдеров WGSL}
\begin{itemize}
\item Есть массивы: \mintinline{wgsl}|let a = array<f32, 5>(0.0, 1.0, 2.0, 3.0, 4.0);|
\pause
\item Константы (известные на момент компиляции): \mintinline{wgsl}|const float PI = 3.141592;|
\end{itemize}
\end{frame}

\begin{frame}[fragile]
\frametitle{Язык описания шейдеров WGSL}
\begin{itemize}
\item Ветвление: \mintinline{wgsl}|if (condition) { ... } else { ... }|
\pause
\item Тернарный if: \mintinline{wgsl}|select(if_false, if_true, condition);|
\pause
\item Циклы: \mintinline{wgsl}|for (var i = 0; i < 10; ++i) { ... }|
\begin{itemize}
\item Бесконечные циклы могут проявиться по-разному, вплоть до креша всей системы :)
\end{itemize}
\pause
\usemintedstyle{lightbulb}
\item Функции:
\begin{tcblisting}{listing engine=minted, minted language=wgsl, minted options={fontsize=\scriptsize}, listing only, colback=codebg, colframe=codebg, rounded corners}
fn diffuse(n: vec3f, l: vec3f) -> f32 {
    return max(0.0, dot(n, l));
}
\end{tcblisting}
\pause
\begin{itemize}
\item Могут вызывать другие функции
\pause
\item Рекурсия \textbf{запрещена}
\end{itemize}
\end{itemize}
\end{frame}

\begin{frame}[fragile]
\frametitle{Язык описания шейдеров WGSL}
\begin{itemize}
\item Про условные конструкции в шейдерах ходит очень много мифов, вот TL;DR:
\pause
\item Старые GPU действительно плохо умели их обрабатывать из-за \textit{wave divergence}: разным ядрам одного wavefront'а приходится выполнять разные операции
\pause
\item Современные GPU хорошо обрабатывают условия, если их результат одинаковый в рамках одного \textit{wavefront'а} (близкие по номеру вершины, близкие на экране пиксели)
\pause
\item Например, когда условие зависит только от uniform или immediate переменных (т.н. \textit{uniform control flow}), его результат один для всех вызовов шейдера (в рамках одной команды рисования)
\end{itemize}
\end{frame}

\begin{frame}[fragile]
\frametitle{Полезные ресурсы о шейдерах}
\begin{itemize}
\item Туториал: \href{https://google.github.io/tour-of-wgsl}{\nolinkurl{google.github.io/tour-of-wgsl}}
\item Туториал: \href{https://webgpufundamentals.org/webgpu/lessons/webgpu-wgsl.html}{\nolinkurl{webgpufundamentals.org/webgpu/lessons/webgpu-wgsl.html}}
\item Книжка-учебник с большим количеством примеров сложных шейдеров: \href{https://thebookofshaders.com/00/}{\texttt{The Book of Shaders}}
\item Онлайн-редактор шейдеров: \href{https://shadertoy.com}{\nolinkurl{shadertoy.com}}
\end{itemize}
\end{frame}

\begin{frame}[fragile]
\frametitle{Четырёхмерные векторы?}
\usemintedstyle{solarized-light}
\begin{itemize}
\item WGSL поддерживает операции с матрицами и векторами размерностей вплоть до 4
\item \mintinline{wgsl}|@builtin(position)| имеет 4 компоненты
\pause
\item Мы хотим рисовать двухмерные и трёхмерные объекты, зачем четвёртая координата?
\pause
\item Три причины:
\pause
\begin{itemize}
\item Для некоторых \textit{типов данных}: кватернионы, цвета (rgba)
\pause
\item Для \textit{перспективной проекции} (4ая лекция)
\pause
\item Для \textit{аффинных преобразований}
\end{itemize}
\end{itemize}
\end{frame}

\begin{frame}[fragile]
\frametitle{Аффинное пространство}
\begin{itemize}
\item Векторные (линейные) пространства -- круто: большая теория, эффективные операции
\pause
\item Не очень подходит для того, чтобы моделировать пространство
\pause
\begin{itemize}
\item Где в пространстве точка, соответствующая нулевому вектору?
\pause
\item Что такое сумма двух точек?
\end{itemize}
\pause
\item Мысль: точки \begin{math}\neq\end{math} векторы
\end{itemize}
\end{frame}

\begin{frame}[fragile]
\frametitle{Аффинное пространство}
\begin{itemize}
\item Аффинное пространство: векторное пространство, в котором ``забыли''  нулевой вектор
\pause
\item Формальнее, аффинное пространство \begin{math}\mathbb A\end{math} над векторным пространством \begin{math}\mathbb V\end{math} -- это:
\pause
\begin{itemize}
\item Множество точек \begin{math}\mathbb A\end{math}
\pause
\item Операция ``точка + вектор'': \begin{math}\oplus: \mathbb A \times \mathbb V \rightarrow \mathbb A\end{math}
\pause
\item Операция ``точка - точка'': \begin{math}\ominus: \mathbb A \times \mathbb A \rightarrow \mathbb V\end{math}
\pause
\item Аксиомы, связывающие их с операциями векторного пространства, например
\begin{math}(p\oplus v_1)\oplus v_2 = p\oplus(v_1+v_2)\end{math}
\end{itemize}
\pause
\item Любую точку \begin{math}p_0\in \mathbb A\end{math} можно принять за начало координат
\item Это даёт биекцию \begin{math}\mathbb A \rightarrow \mathbb V\end{math} по формуле \begin{math}p \mapsto p - p_0\end{math}
\end{itemize}
\end{frame}

\begin{frame}[fragile]
\frametitle{Аффинное пространство}
\begin{itemize}
\item Множество решений уравнения \begin{math}Lx=0\end{math} в векторном пространстве -- \textit{векторное подпространство} \begin{math}\ker L\end{math}
\pause
\item Множество решений уравнения \begin{math}Lx=y\end{math} в векторном пространстве -- \alert{\textbf{не}} векторное подпространство
\pause
\item Множество решений уравнения \begin{math}Lx=y\end{math} в векторном пространстве -- \textit{аффинное пространство} над \begin{math}\ker L\end{math}
\end{itemize}
\end{frame}

\begin{frame}[fragile]
\frametitle{Аффинные комбинации}
\begin{itemize}
\item Пусть есть набор точек \begin{math}p_1, \dots, p_N\end{math} и набор коэффициентов \begin{math}\lambda_1, \dots, \lambda_N\end{math}
\pause
\item Выберем начало координат \begin{math}p_0\end{math} и определим

\begin{math}\sum \lambda_i p_i = p_0 + \sum \lambda_i (p_i - p_0) \in \mathbb A\end{math}
\pause
\item Теорема: если \begin{math}\sum \lambda_i = 1\end{math}, то результат (точка) не зависит от выбора \begin{math}p_0\end{math}
\pause
\item В этом случае \begin{math}q = \sum \lambda_i p_i\end{math} называется \textit{аффинной комбинацией} точек \begin{math}p_i\end{math} с коэффициентами \begin{math}\lambda_i\end{math}
\pause
\item Коэффициенты \begin{math}\lambda_i\end{math} называются \textit{барицентрическими координатами} точки \begin{math}q\end{math}
\pause
\item Поиск \begin{math}\lambda_i\end{math} по известным \begin{math}q\end{math} и \begin{math}p_i\end{math} сводится к решению линейной системы уравнений (квадратная только если \begin{math}N = D + 1\end{math}, где \begin{math}D\end{math} -- размерность пространства)
\pause
\item Аффинная комбинация пустого множества точек не определена
\end{itemize}
\end{frame}

\begin{frame}[fragile]
\frametitle{Аффинные комбинации: пример}
\slideimage{lerp-base.png}
\end{frame}

\begin{frame}[fragile]
\frametitle{Аффинные комбинации: пример}
\slideimage{lerp-1.png}
\end{frame}

\begin{frame}[fragile]
\frametitle{Аффинные комбинации: пример}
\slideimage{lerp-2.png}
\end{frame}

\begin{frame}[fragile]
\frametitle{Аффинные комбинации: пример}
\slideimage{lerp-3.png}
\end{frame}

\begin{frame}[fragile]
\frametitle{Аффинные комбинации: пример}
\slideimage{lerp-4.png}
\end{frame}

\begin{frame}[fragile]
\frametitle{Векторнозначные комбинации}
\begin{itemize}
\item Пусть есть набор точек \begin{math}p_1, \dots, p_N\end{math} и набор коэффициентов \begin{math}\lambda_1, \dots, \lambda_N\end{math}
\pause
\item Выберем начало координат \begin{math}p_0\end{math} и определим

\begin{math}\sum \lambda_i p_i = \sum \lambda_i (p_i - p_0) \in \mathbb V\end{math}
\pause
\item Теорема: если \begin{math}\sum \lambda_i = 0\end{math}, то результат (вектор) не зависит от выбора \begin{math}p_0\end{math}
\end{itemize}
\end{frame}

\begin{frame}[fragile]
\frametitle{Аффинная оболочка}
\begin{itemize}
\item Пусть есть набор точек \begin{math}p_1, \dots, p_N\end{math}
\pause
\item Их \textit{аффинная оболочка} -- множество всех возможных аффинных комбинаций, т.е. множество \begin{math}\sum \lambda_i p_i\end{math} для всех возможных \begin{math}\lambda_i\end{math} таких, что \begin{math}\sum \lambda_i = 1\end{math}
\pause
\item Аффинная оболочка точки -- сама \textit{точка}
\pause
\item Аффинная оболочка двух точек -- содержащая их \textit{прямая}
\pause
\item Аффинная оболочка трёх точек -- содержащая их \textit{плоскость}
\pause
\item И т.д.
\end{itemize}
\end{frame}

\begin{frame}[fragile]
\frametitle{Линейная интерполяция}
\begin{itemize}
\item Линейная интерполяция = интерполяция, использующая барицентрические координаты как коэффициенты
\pause
\item Пусть есть точки \begin{math}p_i\end{math} и точка \begin{math}q = \sum \lambda_i p_i\end{math}
\pause
\item Пусть в точках \begin{math}p_i\end{math} задано значение некоторой величины \begin{math}f_i\end{math}
\pause
\item Линейная интерполяция величины \begin{math}f\end{math} в точке \begin{math}q\end{math} -- значение \begin{math}\sum \lambda_i f_i\end{math}
\pause
\item Для этого нужно вычисляить \begin{math}\lambda_i\end{math} по известным \begin{math}p_i\end{math} и \begin{math}q\end{math}, что сводится к системе линейных уравнений
\pause
\item \textbf{Ровно это} происходит при интерполяции значений, переданных из вершинного шейдера во фрагментный \pause (с точностью до перспективной проекции, о чём мы поговорим позже)
\end{itemize}
\end{frame}

\begin{frame}[fragile]
\frametitle{Линейная интерполяция цветов}
\slideimage{triangle.png}
\end{frame}

\begin{frame}[fragile]
\frametitle{Выпуклые оболочки}
\begin{itemize}
\item Выпуклая комбинация: аффинная комбинация, в которой все коэффициенты неотрицательны \begin{math}\lambda_i \geq 0\end{math}
\pause
\begin{itemize}
\item Вместе с \begin{math}\sum\lambda_i = 1\end{math} это означает, что \begin{math}\lambda_i \in [0, 1]\end{math}
\end{itemize}
\pause
\item Множество всех выпуклых комбинаций набора точек = \textit{выпуклая оболочка}
\pause
\item Выпуклая оболочка точки -- сама \textit{точка}
\pause
\item Выпуклая оболочка двух точек -- \textit{отрезок} между этими точками
\pause
\item Выпуклая оболочка трёх точек -- \textit{треугольник} с этими точками в качестве вершин
\pause
\item Выпуклая оболочка четырёх точек -- \textit{тетраэдр} с этими точками в качестве вершин
\pause
\item И т.д.
\end{itemize}
\end{frame}

\begin{frame}[fragile]
\frametitle{Выпуклые оболочки}
\usemintedstyle{solarized-light}
\begin{itemize}
\item При растеризации центры пикселей попадают внутрь треугольника
\pause
\item \begin{math}\Rightarrow\end{math} Пиксели лежат в \textit{выпуклой оболочке} исходных вершин
\pause
\item \begin{math}\Rightarrow\end{math} Коэффициенты интерполяции всегда в диапазоне \begin{math}\lambda_i \in [0, 1]\end{math}
\pause
\item \begin{math}\Rightarrow\end{math} Проинтерполированные значения лежат в выпуклой оболочке значений в вершинах
\begin{itemize}
\item Например, если \mintinline{cpp}|out|-значения величины в вершинах были \begin{math}1.5\end{math}, \begin{math}0.3\end{math} и \begin{math}5.7\end{math}, то \mintinline{cpp}|in|-значения в пикселях будут в диапазоне \begin{math}[0.3, 5.7]\end{math} \textit{(с точностью до floating-point ошибок)}
\end{itemize}
\pause
\item N.B. Это может нарушаться при \textit{мультисэмплинге}, о котором мы поговорим позднее
\end{itemize}
\end{frame}

\begin{frame}[fragile]
\frametitle{Аффинные преобразования}
\begin{itemize}
\item Линейные преобразования -- сохраняют \textit{линейные комбинации} векторов: \begin{math}S \left(\sum \lambda_i v_i\right) = \sum \lambda_i S v_i\end{math}
\pause
\item Аффинные преобразования -- сохраняют \textit{аффинные комбинации} точек: \begin{math}S \left(\sum \lambda_i p_i\right) = \sum \lambda_i S p_i\end{math}
\pause
\item Если выбрано начало координат \begin{math}p_0\end{math}, преобразование \begin{math}S: \mathbb A \rightarrow \mathbb A\end{math} можно понимать как преобразование соответствующего векторного пространства
\begin{math}S: \mathbb V \rightarrow \mathbb V\end{math} по формуле \begin{math}S(v) = S(p_0 + v) - p_0\end{math}
\pause
\item Можно показать, что в таком виде любое аффинное преобразование это линейное преобразование + сдвиг на фиксированный вектор:

\begin{math}S(v) = Av + b\end{math}
где
\begin{itemize}
\item \begin{math}A\end{math} -- линейное преобразование пространства \begin{math}\mathbb V\end{math}
\item \begin{math}b\end{math} -- вектор из пространства \begin{math}\mathbb V\end{math}
\end{itemize}
\end{itemize}
\end{frame}

\begin{frame}[fragile]
\frametitle{Аффинные преобразования}
\begin{itemize}
\item В коде аффинное преобразование пространства размерности \begin{math}N\end{math} можно представлять как \textit{пару} \begin{math}(A,b)\end{math} из матрицы \begin{math}N\times N\end{math} и вектора размерности \begin{math}N\end{math}
\pause
\item Пара \begin{math}(A,b)\end{math} действует на точку/вектор \begin{math}v\end{math} по формуле

\begin{math}(A, b) \cdot v = Av + b\end{math}
\pause
\item Композиция преобразований:

\begin{math}(A_1, b_1) \cdot (A_2, b_2) \cdot v = (A_1, b_1) \cdot (A_2 v + b_2) = A_1 (A_2 v + b_2) + b_1 = (A_1 A_2) v + (A_1 b_2 + b_1) = (A_1 A_2, A_1 b_2 + b_1) \cdot v\end{math}
\pause
\item Тождественное преобразование:
\begin{math}(\mathbb I, 0)\end{math}
\pause
\item Обратное преобразование:

\begin{math}(A, b)^{-1} = (A^{-1}, -A^{-1}b)\end{math}
\end{itemize}
\end{frame}

\begin{frame}[fragile]
\frametitle{Аффинные преобразования}
\begin{itemize}
\item Повсеместно используются для задания положения, ориентации и размера объектов
\pause
\item Векторная часть преобразования: положение 3D-объекта
\pause
\item Матричная часть преобразования: вращение и размер объекта
\pause
\item При чём тут четвёртая координата?
\end{itemize}
\end{frame}

\begin{frame}[fragile]
\frametitle{Однородные координаты}
\begin{itemize}
\item Трюк: можно вложить \textit{вектор} размерности \begin{math}N\end{math} в векторное пространство размерности \begin{math}N+1\end{math}, добавив \begin{math}0\end{math} в качестве последней координаты
\pause
\item Трюк: можно вложить \textit{точку} размерности \begin{math}N\end{math} в векторное пространство размерности \begin{math}N+1\end{math}, добавив \begin{math}1\end{math} в качестве последней координаты
\pause
\item Такое представление точек и векторов называется \textit{однородными координатами}
\pause
\item Согласуется с операциями на векторах и точках
\pause
\item Согласуется с аффинными комбинациями
\end{itemize}
\end{frame}

\begin{frame}[fragile]
\frametitle{Однородные координаты}
\begin{itemize}
\item Позволяет представить аффинное преобразование \begin{math}N\end{math}-мерного пространства как матрицу \begin{math}(N+1)\times (N+1)\end{math}:

\begin{center}
\begin{math}
\begin{pmatrix}
A & b \\
0 & 1
\end{pmatrix}
\end{math}
\end{center}
\pause
\item Можно применять как к точкам, так и к векторам (векторы игнорируют сдвиг на \begin{math}b\end{math})
\pause
\item Композиция преобразований = умножение матриц
\pause
\item \begin{math}\Rightarrow\end{math} Позволяет удобно комбинировать преобразования (напр. масштабирование + сдвиг + поворот + другой сдвиг)
\pause
\item \begin{math}\Rightarrow\end{math} Позволяет собрать положение и ориентацию камеры и её проекцию в одну матрицу (обсудим на 4ой лекции)
\end{itemize}
\end{frame}

\begin{frame}[fragile]
\frametitle{Точка в однородных координатах}
\begin{center}
\begin{math}
\begin{pmatrix}
p_1 \\
p_2 \\
p_3 \\
1
\end{pmatrix}
\end{math}
\end{center}
\end{frame}

\begin{frame}[fragile]
\frametitle{Вектор в однородных координатах}
\begin{center}
\begin{math}
\begin{pmatrix}
v_1 \\
v_2 \\
v_3 \\
0
\end{pmatrix}
\end{math}
\end{center}
\end{frame}

\begin{frame}[fragile]
\frametitle{Аффинное преобразование в однородных координатах}
\begin{center}
\begin{math}
\begin{pmatrix}
A_{1,1} & A_{1,2} & A_{1,3} & b_1 \\
A_{2,1} & A_{2,2} & A_{2,3} & b_2 \\
A_{3,1} & A_{3,2} & A_{3,3} & b_3 \\
0 & 0 & 0 & 1
\end{pmatrix}
\end{math}
\end{center}
\end{frame}

\begin{frame}[fragile]
\frametitle{Матрица сдвига на фиксированный вектор}
\begin{center}
\begin{math}
\begin{pmatrix}
1 & 0 & 0 & t_x \\
0 & 1 & 0 & t_y \\
0 & 0 & 1 & t_z \\
0 & 0 & 0 & 1
\end{pmatrix}
\cdot
\begin{pmatrix}
p_x \\
p_y \\
p_z \\
1
\end{pmatrix}
=
\begin{pmatrix}
p_x + t_x \\
p_y + t_y\\
p_z + t_z\\
1
\end{pmatrix}
\end{math}

\begin{math}
\begin{pmatrix}
1 & 0 & 0 & t_x \\
0 & 1 & 0 & t_y \\
0 & 0 & 1 & t_z \\
0 & 0 & 0 & 1
\end{pmatrix}
\cdot
\begin{pmatrix}
v_x \\
v_y \\
v_z \\
0
\end{pmatrix}
=
\begin{pmatrix}
v_x\\
v_y\\
v_z\\
0
\end{pmatrix}
\end{math}
\end{center}
\end{frame}

\begin{frame}[fragile]
\frametitle{Матрица изотропного масштабирования}
\begin{center}
\begin{math}
\begin{pmatrix}
s & 0 & 0 & 0 \\
0 & s & 0 & 0 \\
0 & 0 & s & 0 \\
0 & 0 & 0 & 1
\end{pmatrix}
\cdot
\begin{pmatrix}
p_x \\
p_y \\
p_z \\
1
\end{pmatrix}
=
\begin{pmatrix}
s p_x\\
s p_y\\
s p_z\\
1
\end{pmatrix}
\end{math}

\begin{math}
\begin{pmatrix}
s & 0 & 0 & 0 \\
0 & s & 0 & 0 \\
0 & 0 & s & 0 \\
0 & 0 & 0 & 1
\end{pmatrix}
\cdot
\begin{pmatrix}
v_x \\
v_y \\
v_z \\
0
\end{pmatrix}
=
\begin{pmatrix}
s v_x\\
s v_y\\
s v_z\\
0
\end{pmatrix}
\end{math}
\end{center}
\end{frame}

\begin{frame}[fragile]
\frametitle{Матрица поворота на угол \begin{math}\theta\end{math} в плоскости XY}
\begin{center}
\begin{math}
\begin{pmatrix}
\cos \theta & - \sin \theta & 0 & 0 \\
\sin \theta & \cos \theta & 0 & 0 \\
0 & 0 & 1 & 0 \\
0 & 0 & 0 & 1
\end{pmatrix}
\cdot
\begin{pmatrix}
p_x \\
p_y \\
p_z \\
1
\end{pmatrix}
=
\begin{pmatrix}
p_x \cos \theta - p_y \sin \theta \\
p_x \sin \theta + p_y \cos \theta \\
p_z \\
1
\end{pmatrix}
\end{math}

\begin{math}
\begin{pmatrix}
\cos \theta & - \sin \theta & 0 & 0 \\
\sin \theta & \cos \theta & 0 & 0 \\
0 & 0 & 1 & 0 \\
0 & 0 & 0 & 1
\end{pmatrix}
\cdot
\begin{pmatrix}
v_x \\
v_y \\
v_z \\
0
\end{pmatrix}
=
\begin{pmatrix}
v_x \cos \theta - v_y \sin \theta \\
v_x \sin \theta + v_y \cos \theta \\
v_z \\
0
\end{pmatrix}
\end{math}
\end{center}
\end{frame}

\begin{frame}
\frametitle{Аффинные преобразования: ссылки}
\begin{itemize}
\item \href{https://learnopengl.com/Getting-started/Transformations}{\nolinkurl{learnopengl.com/Getting-started/Transformations}}
\item \href{https://open.gl/transformations}{\nolinkurl{open.gl/transformations}}
\item \href{https://en.wikipedia.org/wiki/Affine_space}{\nolinkurl{en.wikipedia.org/wiki/Affine\_space}}
\item \href{https://en.wikipedia.org/wiki/Affine_transformation}{\nolinkurl{en.wikipedia.org/wiki/Affine\_transformation}}
\end{itemize}
\end{frame}

\end{document}